\documentclass[12pt]{article}
\usepackage{bbm}
\usepackage{fancyhdr}
\usepackage{titling}
\usepackage{enumitem}
\usepackage{amsmath}
\usepackage{dsfont}
\usepackage{amsfonts}
\usepackage{amssymb}
\usepackage{amsthm}
\usepackage{geometry}
\usepackage{tikz}
\usepackage{algorithm}
\usepackage{algpseudocode}
\geometry{verbose,tmargin=1in,bmargin=1in,lmargin=1in,rmargin=1in}
\renewcommand{\vec}[1]{\mathbf{#1}}
\newcommand{\dec}{\text{dec}}
\newcommand{\divi}{\mathbin{\textnormal{div}}}
\title{Homework 1}
\author{Satvik Prasad}
\date{}
\newtheorem{lemma}{Lemma}[section]
\fancypagestyle{plain}{%
  \fancyhf{}
  \fancyhead[L]{\theauthor}
  \fancyhead[R]{\thedate}
}
\pagestyle{plain}
\makeatletter
\def\@maketitle{%
  \newpage
  \null
  \vskip 1em%
  \begin{center}%
    \let \footnote \thanks
    {\LARGE \@title \par}%
    \vskip 1em%
  \end{center}%
  \par
\vskip 1em}
\makeatother
\usepackage{cmbright}

\newcommand{\mat}[1]{\mathbf{#1}}

% ---------- Probability theory ----------
% Core operators (optional argument-free; use as \E{X}, \Var{X}, \Cov{X}{Y})
\DeclareMathOperator{\EOp}{\mathbb{E}}
\DeclareMathOperator{\VarOp}{Var}
\DeclareMathOperator{\CovOp}{Cov}
\DeclareMathOperator{\CorrOp}{Corr}

\renewcommand{\P}[1]{\mathbb{P}\left(#1\right)}        % \P{A}
\newcommand{\E}[1]{\EOp\left[#1\right]}                % \E{X}
\newcommand{\Var}[1]{\VarOp\left(#1\right)}            % \Var{X}
\newcommand{\Cov}[2]{\CovOp\left(#1,\,#2\right)}       % \Cov{X}{Y}
\newcommand{\Corr}[2]{\CorrOp\left(#1,\,#2\right)}     % \Corr{X}{Y}
\newcommand{\SD}[1]{\operatorname{SD}\left(#1\right)}  % \SD{X}

% Subscripted versions, e.g. \Esub{\theta}{X}, \Psub{\theta}{A}
\newcommand{\Esub}[2]{\EOp_{#1}\left[#2\right]}
\newcommand{\Psub}[2]{\mathbb{P}_{#1}\left(#2\right)}

% Conditioning, indicators, independence
\newcommand{\given}{\,\middle|\,}                      % \E{X \given Y}
\newcommand{\ind}[1]{\mathbbm{1}\left\{#1\right\}}     % \ind{X > 0}
\newcommand{\indep}{\perp\!\!\!\perp}                  % X \indep Y
\newcommand{\iid}{\overset{\text{iid}}{\sim}}          % X_i \iid \Normal(0,1)

% Convergence
\newcommand{\convp}{\xrightarrow{\;p\;}}               % in probability
\newcommand{\convd}{\xrightarrow{\;d\;}}               % in distribution
\newcommand{\convas}{\xrightarrow{\text{a.s.}}}        % almost surely
\newcommand{\convLp}[1]{\xrightarrow{L^{#1}}}          % in L^p

% Probability spaces
\newcommand{\F}{\mathcal{F}}                           % sigma-algebra
\newcommand{\B}{\mathcal{B}}                           % Borel sigma-algebra
\newcommand{\pspace}{(\Omega, \F, \mathbb{P})}

% Common distributions
\newcommand{\Normal}{\mathcal{N}}
\newcommand{\Bern}{\operatorname{Bernoulli}}
\newcommand{\Bin}{\operatorname{Binomial}}
\newcommand{\Pois}{\operatorname{Poisson}}
\newcommand{\Unif}{\operatorname{Uniform}}
\newcommand{\Exp}{\operatorname{Exp}}
\newcommand{\Geom}{\operatorname{Geometric}}
\newcommand{\Gam}{\operatorname{Gamma}}
\newcommand{\Beta}{\operatorname{Beta}}

% Number sets
\newcommand{\R}{\mathbb{R}}
\newcommand{\N}{\mathbb{N}}
\newcommand{\Z}{\mathbb{Z}}

\begin{document}
\begin{enumerate}
    \item \begin{enumerate}
        \item If $\mathbb P(A \mid B) \leq P(A),$ then we know \begin{align*}
            \mathbb P(A \cap B) \leq \mathbb P(A) \mathbb P(B).
        \end{align*}

        Then we know $$\mathbb P(A \mid B^c) = \dfrac{\mathbb P(A \cap B^c)}{1 - \mathbb P(B)}.$$
        Since $(A \cap B) \cup (A \cap B^c) = A,$ we can deduce that \begin{align*}
            \mathbb P(A \mid B^c) = \dfrac{\mathbb P(A) - \mathbb P(A \cap B)}{1 - \mathbb P(B)} \geq \dfrac{\mathbb P(A) - \mathbb P(A) \mathbb P(B)}{1 - \mathbb P(B)} = \mathbb P(A),
        \end{align*}
        since $\mathbb P(B) < 1.$ Therefore, this statement is true.

        \item True, since $\mathbb P(A \cap B) + \mathbb P(A^c \cap B) = \mathbb P(B).$

        \item False.

        \item False. Their intersections with $B$ don't need to have equal probability. 
    \end{enumerate}

    \item We want to find $\mathbb P(X = B \mid Y = B),$ where $Y$ represents the observed color, and $X$ represents the true color. Using Bayes' theorem, we get 
    \begin{align*}
        \mathbb P(X = B \mid Y = B)  = \dfrac{\mathbb P(X = B)}{\mathbb P(Y = B)} \cdot \mathbb P(Y = B \mid X = B) &=  \dfrac{0.15}{0.15 \cdot 0.8 + 0.85 \cdot 0.2} \cdot 0.8 \\
        &\approx 41\%.
    \end{align*}

    \item \begin{enumerate}
        \item Using the definition of conditionality, we get \begin{align*}
            \mathbb P(A \cap B \mid  C) = \dfrac{\mathbb P(A \cap B \cap C)}{\mathbb P(C)} = \dfrac{\mathbb P((A \cap C) \cap (B \cap C))}{\mathbb P(C)}. 
        \end{align*}
    However, $A \cap C$ and $B \cap C$ are not necessarily independent events given $A$ and $B$ are independent.

    For example, $A = \{X_1 = H\}$ and $B = \{X_2 = H\}$ are independent for a fair coin, but conditioned on the two flips being the same face, they are no longer independent since $\mathbb P(A \cap B \mid C) = 1/2,$ while $\mathbb P(A \mid C) \mathbb P(B \mid C) = 1/4.$


    \item In other words, we are given that \begin{align*}
        \mathbb P(A \cap B \cap C) = \mathbb P(A \cap C) \mathbb P(B \cap C).
    \end{align*}

    Once again, consider the case where the probability of our coin being double sided is $10\%.$ Given that it is not double-sided (C), $A$ and $B$ are independent, however in the usual sense $A$ and $B$ are not independent, since there is a chance that the coin is biased.
    \end{enumerate}

    \item For the following questions, let $B$ be the event you have the black plague, and $T$ represent the test outcome. 
    \begin{enumerate}
    \item Using Bayes', \begin{align*}
        \mathbb P(B \mid T) = \dfrac{\underbrace{0.98}_{\mathbb P(T \mid B)} \cdot \underbrace{0.01}_{\mathbb P(B)}}{\underbrace{0.01 \cdot 0.98 + 0.99 \cdot 0.02}_{\mathbb P(T)}} \approx 33\%
    \end{align*}

    \item Using Bayes', \begin{align*}
        \mathbb P(B \mid T) = \dfrac{0.98^2 \cdot 0.01}{0.01 \cdot 0.98^2 + 0.99 \cdot 0.02^2} \approx 96\%.
    \end{align*}
    \end{enumerate}

    \item \begin{enumerate}
        \item At position $k,$ you have $1/2$ probability to move right, or move left. Let $X_k = \{\text{position $0$ before position $n$ starting at $k$}\}.$ Let $H$ represent the current roll being a heads, and $T$ the converse. Hence, since \begin{align*}
            X_k = \underbrace{(X_{k + 1} \cap H) \cup (X_{k -1 } \cap T)}_{\text{disjoint}}.
        \end{align*}
        Thus, \begin{align*}
            p_k = \mathbb P(X_k) = \mathbb P(X_{k + 1} \cap H) + \mathbb P(X_{k-1} \cap T).
        \end{align*}
        $X_j$ depends only on future coin flips, which are independent of the current coin flip, so we can say \begin{align*}
            p_k &= \mathbb P(H) \cdot \mathbb P(X_{k + 1}) + \mathbb P(T) \cdot \mathbb P(X_{k -1}) \\
            &= \dfrac{1}{2} \mathbb P(X_{k + 1}) + \dfrac{1}{2} \mathbb P(X_{k -1}) \\
            &= \dfrac{1}{2} p_{k + 1} + \dfrac{1}{2} p_{k -1 }
        \end{align*}
        as required.

        \item $p_0$ should take value $1,$ since you are already at position $0.$ Similarly, $p_n$ should take value $0,$ since you are already at position $n.$ Then, flattening the recurrence, we get \begin{align*}
            p_{n - 1}  &= \dfrac{1}{2}p_{n - 2}, \\
            p_{n - 2}  &= \dfrac{2}{3}p_{n - 3}, \\
            p_{n - 3} &= \dfrac{3}{4} p_{n - 4},  \\
            &\ldots 
        \end{align*}
        Thus, the sequence \begin{align*}
            \{p_k\}_{0 \leq k \leq n} = \left\{1, \dfrac{n - 1}{n}, \dfrac{n -2}{n}, \ldots, \dfrac{1}{n}, 0\right\}
        \end{align*}
        satisfies this recurrence.
    \end{enumerate}

    \item \begin{enumerate}
        \item The parallel circuit only requires at least one component to operate, meaning the probability is $$1 - \left(\frac{1}{2}\right)^N.$$
        
        \item Requires all of the components to function, yielding $$\left(\frac{1}{2}\right)^N.$$

        \item The probability of the first component functioning and the entire circuit functioning is just the probability of the first component functioning, meaning $$p  =\dfrac{2^{N-1}}{2^N - 1}.$$
        
        \item Let $A = \{\text{first component functioning}\}$, $B = \{\text{entire circuit functioning}\}$, $C=\{\text{entire circuit functioning}\}$ Since $P(A \cap B) = \mathbb P(C),$ we know $$p = 1.$$
    \end{enumerate}

    For large $N,$ (a) goes to $1,$ (b) goes to $0,$ (c) tends to $\dfrac{1}{2},$ and (d) is always $1.$

    \item We know that $\mathbb P(A \cap B) = 1/4 = \mathbb P(A) \mathbb P(B),$ $\mathbb P(B \cap C) = \mathbb P(\text{both even}) = 1/4 = \mathbb P(B) \mathbb P(C),$ and $\mathbb P(A \cap C) = \mathbb P(\text{both odd}) = 1/4 = \mathbb P(A) \mathbb P(C),$ meaning all $A, B, C$ are pairwise independent.

    However, $\mathbb P(A \cap B \cap C) = 0 \neq 1/8 = \mathbb P(A) \mathbb P(B) \mathbb P(C),$ meaning the events aren't all independent.
\end{enumerate}
\end{document}
